% Preamble / 宏包区
\documentclass{article}
\usepackage[margin=0mm]{geometry}
\usepackage{xcolor}
\usepackage{graphicx}
\usepackage{tikz}
\usepackage{amsmath}

\begin{document}
% Metadata / 注释区
% Exported by guiLaTeX ExportCore
% Export date: 2026-04-12 16:02:40
% Page size: A4 (210mm × 297mm)

% Optional semantic summary 区
% Title: 标题：可视化编辑 LaTeX 文档 (文本框)
% Author: 作者：某某 / 博士生原型测试 (文本框)
% Absolute positioned objects 区
\begin{tikzpicture}[remember picture, overlay]
  % Layer 1: 1
  
ode[anchor=north west, rotate=0, text width=400mm, align=left, font=ontsize{24}{28.799999999999997}\selectfont, color={rgb,1:red,0.0;green,0.0;blue,0.0}, visible on layer=1]
    at (100mm, 197mm) {标题：可视化编辑 LaTeX 文档 (文本框)};
  % Layer 2: 2
  
ode[anchor=north west, rotate=0, text width=200mm, align=left, font=ontsize{16}{19.2}\selectfont, color={rgb,1:red,0.0;green,0.0;blue,0.0}, visible on layer=2]
    at (100mm, 147mm) {作者：某某 / 博士生原型测试 (文本框)};
  % Layer 3: 3
  
ode[anchor=north west, rotate=0, text width=400mm, align=left, font=ontsize{14}{16.8}\selectfont, color={rgb,1:red,0.0;green,0.0;blue,0.0}, visible on layer=3]
    at (100mm, 97mm) {这是一段用于演示段落对象的正文内容，可用于测试拖动、缩放与排版差异。(段落)};
  % Layer 4: 4
  
ode[anchor=north west, rotate=0, text width=150mm, align=left, font=ontsize{18}{21.599999999999998}\selectfont, color={rgb,1:red,0.0;green,0.0;blue,0.0}, visible on layer=4]
    at (100mm, 17mm) {E = mc² (文本框)};
  % Layer 5: 5
  
ode[anchor=north west, rotate=0, text width=200mm, align=left, font=ontsize{12}{14.399999999999999}\selectfont, color={rgb,1:red,0.0;green,0.0;blue,0.0}, visible on layer=5]
    at (100mm, -33mm) {\includegraphics[width=200mm, height=150mm]{图片示例 / 占位图}};
\end{tikzpicture}
\end{document}